% ============================================================================
%  SECCIÓN 3.11: CREACIÓN DEL PROYECTO
% ============================================================================

\section{Creaci\'on del proyecto}

La aplicaci\'on m\'ovil se desarroll\'o con el marco de trabajo Expo, basado en
React Native, que permite crear proyectos multiplataforma de forma r\'apida y
con una configuraci\'on m\'inima \citep{expo}. La creaci\'on del proyecto se
realiz\'o mediante la herramienta de l\'inea de comandos \texttt{create-expo-app},
empleando la plantilla \texttt{blank-typescript}, que incluye TypeScript
configurado por defecto.

La creaci\'on del proyecto se realiz\'o con el gestor de paquetes
\texttt{pnpm}, mediante la ejecuci\'on del siguiente comando:

\begin{lstlisting}
pnpm create expo-app lugares-interactivos --template blank-typescript
\end{lstlisting}

El comando descarga la plantilla, crea la estructura del proyecto en el
directorio \texttt{lugares-interactivos} e instala autom\'aticamente las
dependencias b\'asicas, entre las que destacan \texttt{expo},
\texttt{react}, \texttt{react-native} y \texttt{expo-status-bar}. El archivo
\texttt{pnpm-lock.yaml} generado garantiza que las dependencias se instalen de
manera determinista en todos los equipos del equipo de trabajo
\citep{nodejs}.

En la figura se presenta la evidencia de la creaci\'on del proyecto en la
terminal.

\begin{figure}[H]
\centering
  \caption{Evidencia de la creaci\'on del proyecto con Expo}
  \label{fig:creacion-proyecto}
  \vspace{6pt}
  \begin{lstlisting}[style=terminal]
$ pnpm create expo-app lugares-interactivos --template blank-typescript
Creating a new Expo project...
Downloading the selected template: blank-typescript
Scaffolding project in lugares-interactivos/
Installing dependencies using pnpm...
Packages: +1168
Progress: resolved 1168, reused 1168, downloaded 0, added 1168
Prepared expo in 6.2s

[OK] Your project is ready!

To run your project, navigate to the directory and run one of:
  - pnpm start    # start the development server
  - pnpm android  # open the app on an Android device
  \end{lstlisting}
  \par\raggedright\fuenteAPA{Elaboraci\'on propia en base a la ejecuci\'on del proyecto.}\end{figure}

Una vez creado el proyecto, la ejecuci\'on de los comandos de pnpm se realiza
desde el directorio del proyecto, tal como se describe en la secci\'on de
ejecuci\'on de la aplicaci\'on.
